\documentclass[10pt]{beamer}

\usepackage[utf8]{inputenc}
\usepackage[spanish]{babel}
\usepackage{graphicx}

\mode<presentation>
\usetheme{Madrid}
%\usecolortheme[RGB={111,73,135}]{structure}
\usecolortheme[RGB={128,0,0}]{structure}
%\usecolortheme[RGB={0,96,0}]{structure}
%\usecolortheme[RGB={200,0,200}]{structure}
%\usecolortheme[RGB={0,128,0}]{structure}
%\usecolortheme[RGB={0,0,128}]{structure}
\usefonttheme{serif}
\useinnertheme{rectangles}
\useoutertheme{split}
\setbeamercovered{transparent}


% Definiciones para usar luego
\newtheorem{ejemplo}{Ejemplo}
\newtheorem{definicion}{Definición}

\title{Mejorando el rendimiento de tu Linux}
\author{Jesús Espino García}
\date{13 de Marzo de 2009}
\subject{Mejorando el rendimiento de tu Linux}

\institute[GUL UC3M]{
  Grupo de Usuarios de Linux\\
  Universidad Carlos III de Madrid.\\
  \ \\
  \includegraphics[height=1.5cm]{imgs/gul}}

\setcounter{tocdepth}{2}

\AtBeginSubsection[]
{
  \begin{frame}<beamer>{Indice}
    \tableofcontents[sectionstyle=show/shaded,subsectionstyle=show/shaded/hide]
  \end{frame}
}

\begin{document}

  \frame{\maketitle}
  
  \section{Introducción}
\subsection{¿Qué es?}
\begin{frame}{¿Qué es?}
	\begin{itemize}
		\item Adaptar el equipo a nuestras necesidades.
		\item Equilibrar la capacidad y la demanda.
		\item Dimensionar correctamente nuestro hardware.
	\end{itemize}
\end{frame}

\subsection{¿Por qué?}
\begin{frame}{¿Por qué?}
	\begin{itemize}
		\item La configuración genérica no cubre las necesidades.
		\item El rendimiento es una pieza clave del sistema.
		\item No todos los recursos son igual de importantes.
	\end{itemize}
\end{frame}

\subsection{¿Cuándo?}
\begin{frame}{¿Cuándo?}
	\begin{itemize}
		\item Desde la planificación del proyecto.
		\item En cada una de las etapas del proyecto.
		\item Después de un backup!!!
	\end{itemize}
\end{frame}

\subsection{¿Dónde?}
\begin{frame}{¿Dónde?}
	\begin{itemize}
		\item En el hardware.
		\item En el software.
		\item En el entorno.
	\end{itemize}
\end{frame}

\subsection{¿Cómo?}
\begin{frame}{¿Cómo?}
	\begin{itemize}
		\item Fuertes conocimientos sobre el sistema.
		\item Monitorización.
		\item Benchmarks.
		\item Prueba y Error.
		\item Mas arte que ciencia.
		\item Cambiando:
		\begin{itemize}
			\item Parámetros del kernel.
			\item Parámetros del sistema de ficheros.
			\item Parámetros de los drivers.
			\item Ampliando hardware.
			\item Mejorando el entorno.
			\item Modificando el software (kernel incluido).
		\end{itemize}
	\end{itemize}
\end{frame}

  \section{Entendiendo Linux}
\subsection{Gestión de procesos}
\begin{frame}{¿Qué es un proceso?}
	\begin{itemize}
		\item Estructura de datos en el SO.
		\item Proceso completo o hilo.
		\item Instancia de tarea a realizar por la CPU.
		\item Utiliza, reserva y comparte recursos.
	\end{itemize}
\end{frame}

\begin{frame}{Prioridades}
	\begin{itemize}
		\item Los procesos tienen 2 prioridades.
		\item Estática (Asignada por el tipo de tarea y el nice)
		\item Dinámica (Calculada por el scheduler)
		\item El nivel nice determina el tamaño del quamto del proceso.
	\end{itemize}
\end{frame}

\begin{frame}{Context Switch}
	\begin{itemize}
		\item Para simular multitarea se va alternando la ejecución de procesos.
		\item Cada cambio de contexto tiene unos costes de tiempo.
	\end{itemize}
\end{frame}

\begin{frame}{Interrupciones}
	\begin{itemize}
		\item El hardware y el software produce interrupciones.
		\item Las interrupciones, normalmente producen cambios de contexto.
		\item Las interrupciones hardware, por lo general, se envían a todos los procesadores.
	\end{itemize}
\end{frame}

\begin{frame}{Scheduler}
	\begin{itemize}
		\item Es el encargado de seleccionar un proceso u otro.
		\item Antes del 2.6.23 el scheduler era el O(1).
		\item En el 2.6.23 y posteriores el scheduler es CFS.
	\end{itemize}
\end{frame}

\subsection{Gestión de memoria}
\begin{frame}{Memoria principal}
	\begin{itemize}
		\item Segmentos: Linux apenas los utiliza.
		\item Paginas: Lo que utiliza normalmente Linux.
		\item En sistemas de 32bits se debe dividir la memoria en 2 partes:
		\begin{itemize}
			\item Una accesible directamente (hasta los 896M).
			\item Y el resto mapeada a través de la zona accesible directamente (hasta 64G usando PAE).
		\end{itemize}
	\end{itemize}
\end{frame}

\begin{frame}{Sistema de caches}
	\begin{itemize}
		\item Los sistemas modernos tienen caches de memoria.
		\item Las caches son memorias mas rápidas que evitan el acceso a la memoria principal
		\item Usan ciertas políticas y algoritmos para determinar de que se guarda cache.
	\end{itemize}
\end{frame}

\begin{frame}{Swap}
	\begin{itemize}
		\item Memoria basada en dispositivos de disco.
		\item Almacena paginas de memoria a las que no se acceden en mucho tiempo.
		\item En caso de falta de falta de memoria principal se utiliza la swap como tal.
	\end{itemize}
\end{frame}

\begin{frame}{Asignación de Memoria}
	\begin{itemize}
		\item Al solicitar memoria, si hay paginas libres las asigna.
		\item Si no hay paginas libres, utiliza métodos de reclamado de pagina.
		\item Para esto utiliza el Buddy System, que se encarga de:
		\begin{itemize}
			\item Mantener una lista de las paginas libres.
			\item Solicita la liberación paginas a kswapd.
			\item Evitar la fragmentación de la memoria.
		\end{itemize}
		\item kswapd libera paginas basándose en ciertas políticas.
	\end{itemize}
\end{frame}

\subsection{Gestión de IO}
\begin{frame}{Sistemas de ficheros}
	\begin{itemize}
		\item VFS: Capa de abstracción para el acceso a sistemas de ficheros.
		\item Journaling: Log de las operaciones realizadas.
		\item Sistema de Ficheros: Estructura de datos para almacenar ficheros (normalmente en disco).
	\end{itemize}
\end{frame}

\begin{frame}{Acceso a disco (IO)}
	\begin{itemize}
		\item Los sistemas de ficheros acceden a disco a través del sistema IO.
		\item El sistema IO organiza los acceso a discos con el IO Scheduler.
		\item Y este finalmente escribe los datos a través del driver de disco.
	\end{itemize}
\end{frame}

\begin{frame}{IO Schedulers}
	\begin{itemize}
		\item Se pueden escoger entre 4:
		\begin{itemize}
			\item Anticipatory: Intenta guardar datos para escribir en tramas largas y continuas. Además usa deadline.
			\item CFQ (Completely Fair Queuing): Implementa un reparto equitativo del IO entre procesos.
			\item Deadline: Round robin con deadline.
			\item NOOP: No hace ninguna operación de orden. Es un FIFO.
		\end{itemize}
	\end{itemize}
\end{frame}

\subsection{Métricas de rendimiento}
\begin{frame}{Métricas de procesador}
	\begin{itemize}
		\item Uso de CPU.
		\item User time: Tiempo utilizado por los procesos de usuario (no kernel).
		\item System time: Tiempo utilizado por operaciones del kernel.
		\item Waiting: Tiempo esperando al subsistema IO.
		\item Idle time: Tiempo esperando tareas.
		\item Nice time: Tiempo utilizado para restablecer prioridades de los procesos.
		\item Load average: Carga del equipo (no es un porcentaje).
		\item Runable processes: Numero de procesos listos para correr.
		\item Blocked: Procesos esperando IO.
		\item Context Switch: Cantidad de cambios de contexto.
		\item Interrupts: Cantidad de interrupciones
	\end{itemize}
\end{frame}

\begin{frame}{Metricas de memoria}
	\begin{itemize}
		\item Memoria libre.
		\item Uso de Swap.
		\item Buffer y cache: Caches asignadas al sistema de ficheros o a dispositivos de bloque. 
		\item Slabs: Uso de memoria del kernel (no swapeable).
		\item Memoria activa/inactiva: La memoria inactiva es la que es candidata de ser swapeada.
	\end{itemize}
\end{frame}

\begin{frame}{Metricas de red}
	\begin{itemize}
		\item Paquetes recibidos y enviados.
		\item Bytes recibidos y enviados.
		\item Colisiones por segundo.
		\item Paquetes "tirados".
		\item Overruns: Veces que la tarjeta de red se ha quedado sin espacio en el buffer.
		\item Errors: Numero de tramas marcadas como fallidas.
	\end{itemize}
\end{frame}

\begin{frame}{Metricas de IO}
	\begin{itemize}
		\item iowait: Tiempo de CPU gastado esperando al subsistema IO.
		\item Media de longitud de la cola: Media de peticiones IO en cola.
		\item Media de espera: Tiempo medio de respuesta de una petición IO.
		\item Transferencias por segundo: Numero de operaciones IO por segundo.
		\item Bloques por segundo r/w
		\item KB por segundo r/w
	\end{itemize}
\end{frame}

  \section{Tunning}
\subsection{Herramientas}

\begin{frame}{Monitorizacion}
	\begin{itemize}
		\item top
		\item vmstat
		\item uptime
		\item ps
		\item free
		\item iostat
		\item mpstat
		\item numastat
		\item pmap
		\item netstat
		\item wireshark
		\item strace
		\item /proc
	\end{itemize}
\end{frame}

\begin{frame}{Benchmarks}
	\begin{itemize}
		\item lmbench
		\item iozone
		\item netperf
	\end{itemize}
\end{frame}

\subsection{En la CPU}
\begin{frame}{En la CPU}
	\begin{itemize}
		\item Herramientas: uptime, top, vmstat, lmbench\dots
		\item Determinar prioridades: Baja latencia, alta cantidad de proceso.
		\item Determinar las prioridades de los procesos.
		\item Buscar una disminución del context-switch por interrupción.
		\item Asignar interrupciones a una CPU en concreto (por ejemplo las de una tarjeta red).
		\item Podemos ampliar CPU.
	\end{itemize}
\end{frame}

\subsection{En la memoria}
\begin{frame}{En la memoria}
	\begin{itemize}
		\item Herramientas: top, vmstat, free, pmap, iostat, lmbench\dots
		\item Determinar prioridades: Aprovechamiento de la memoria, menor carga de proceso.
		\item Si hay mucho uso de swap, normalmente indica falta de memoria.
		\item Si hay muchas peticiones de paginas nos podemos plantear ampliar el tamaño de las mismas.
		\item Si tenemos poca memoria y muchos procesos muy pequeños, podemos disminuir el tamaño de pagina. 
		\item Podemos analizar procesos concretos para ver en que consumen la memoria.
		\item Podemos asegurarnos de usar programas con bibliotecas dinámicas.
		\item Podemos ampliar memoria.
	\end{itemize}
\end{frame}

\subsection{En el disco}
\begin{frame}{En la disco}
	\begin{itemize}
		\item Herramientas: vmstat, iostat, iozone, lmbench\dots
		\item Determinar prioridades: Aprovechamiento, menor carga de proceso, mas ancho de banda.
		\item Deberíamos determinar el sistema de fichero que nos interesa.
		\item Si usamos solo ficheros muy grandes, podríamos ampliar el tamaño de bloque del sistema de ficheros.
		\item Si usamos muchos ficheros pequeños, podríamos disminuir el tamaño de bloque del sistema de ficheros.
		\item Escoger el IO Scheduler que nos pueda interesar.
		\item Añadir mas RAM. Esto amplia la memoria de cache para disco.
		\item Establecer parámetros del sistema de ficheros (por ejemplo, noatime).
		\item Establecer las políticas de journaling correctas.
		\item Uso de sistemas RAID o volúmenes lógicos.
	\end{itemize}
\end{frame}

\subsection{En la red}
\begin{frame}{En la red}
	\begin{itemize}
		\item Herramientas: netperf, lmbench, netstat \dots
		\item Determinar prioridades: Ancho de banda, baja latencia.
		\item Ampliar el tamaño del frame (jumbo frames, no soportados por todo el hardware!).
		\item Asignar interrupciones a un procesador en concreto.
		\item Usar mejores tarjetas de red (para mas ancho de banda o para mas baja latencia).
		\item Hacer bonding o balanceo de carga.
		\item El firewall ralentiza las comunicaciones.
	\end{itemize}
\end{frame}

  \section{Para terminar.}

\begin{frame}[containsverbatim]{Referencias}
  \begin{itemize}
    \item http://www.redbooks.ibm.com/redpapers/pdfs/redp4285.pdf: RedPaper de IBM.
  \end{itemize}
\end{frame}

\begin{frame}{Dudas}
  \dots
\end{frame}

\begin{frame}
  \frametitle{\begin{center}Fin\end{center}}
  \begin{center}
    \includegraphics[height=3cm]{imgs/gul}
  \end{center}
\end{frame}


\end{document}
